\documentclass[conference,a4paper]{IEEEtran}
\IEEEoverridecommandlockouts
\usepackage{cite}
\usepackage{amsmath,amssymb,amsfonts}
\usepackage{graphicx}
\usepackage{textcomp}
\usepackage{xcolor}
\usepackage{booktabs}
\usepackage{url}

\def\BibTeX{{\rm B\kern-.05em{\sc i\kern-.025em b}\kern-.08em
    T\kern-.1667em\lower.7ex\hbox{E}\kern-.125emX}}

\begin{document}

\title{Evaluating Squeeze-and-Excitation YOLOv8 for\\
Efficient Online Gambling Content Detection}

\author{\IEEEauthorblockN{Anonymous Author(s)}
\IEEEauthorblockA{Affiliation and contact details withheld\\
for double-blind review}}

\maketitle

\begin{abstract}
Adding architectural complexity to object detectors often raises accuracy.
Its benefit on small, domain-specific datasets is not guaranteed. This
paper studies how three modifications to YOLOv8n trade accuracy against
model complexity. We compare a squeeze-and-excitation (SE) channel-attention
variant, a heavier Res2Net variant with space-to-depth downsampling, and a
lightweight ShuffleNet variant. All models are evaluated against the
unmodified baseline. We use five-fold cross-validation on two datasets. The
first is a small five-class gambling-logo set (Judol Detection v2, 443 images). The second
is a larger two-class Instagram-reel set (Instagram Reels, 3000 images). On the
discriminative dataset, the SE variant reaches the best mean mAP50-95. It
adds only 0.36\% more parameters. The heavier Res2Net variant adds 40.9\%
parameters without any accuracy gain. The lightweight variant trades
accuracy for a 14.5\% parameter reduction. The results suggest that
near-zero-cost channel recalibration is a reliable operating point under
limited data. We report per-fold variability and make no claim of
statistical significance, given the small fold count.
\end{abstract}

\begin{IEEEkeywords}
object detection, YOLOv8, channel attention, squeeze-and-excitation,
model efficiency, small datasets, gambling-logo detection
\end{IEEEkeywords}

\section{Introduction}
Online gambling promotion has become a large moderation problem. Operators
spread brand logos across social media and streaming platforms. These promos
often target young audiences. In many jurisdictions the activity is illegal.
Manual review does not scale to the volume of uploaded content. Automated
visual detection is therefore attractive. A detector must find small gambling
logos inside cluttered frames. It must also run fast enough for high-throughput
moderation. This sets up a classic accuracy versus efficiency problem.

Single-stage detectors are popular for such applied vision tasks. The YOLO
family combines competitive accuracy with real-time speed~\cite{redmon2016you,
bochkovskiy2020yolov4, jocher2023yolov8}. Practitioners often adapt these
detectors to a narrow domain. A common instinct is to enlarge the network to
raise accuracy.

On large benchmarks such as COCO, added capacity usually
helps~\cite{lin2014microsoft}. On small curated datasets, extra capacity can
be unnecessary or even harmful. It also carries a real deployment cost. The
trade-off therefore deserves direct study.

This paper examines that trade-off. We do not chase a large accuracy jump.
Instead we ask which modification gives the best accuracy per unit of added
complexity. We fix YOLOv8n as a common baseline~\cite{jocher2023yolov8}. We
then study three modifications at different points on the complexity axis.
The first is a squeeze-and-excitation block~\cite{hu2018squeeze}. It
recalibrates feature channels at almost no cost. The second combines a
Res2Net multi-scale block~\cite{gao2021res2net} with space-to-depth
downsampling~\cite{sunkara2022no}. It enlarges the model substantially. The
third is a ShuffleNet-based~\cite{ma2018shufflenet} variant with attention.
It reduces the model below the baseline.

Our contribution is empirical, not a new state of the art. We make three
points. First, we evaluate all four models under identical five-fold
cross-validation on two datasets of contrasting size. Second, we report
accuracy jointly with parameter count and computation. Each modification is
judged by cost as well as return. Third, we show that the near-zero-cost SE
variant gives the best mean accuracy on the small dataset. Structural
enlargement does not. This supports an efficiency-first design choice under
limited data.

\section{Related Work}

\subsection{Single-Stage Object Detection}
Two-stage detectors such as Faster R-CNN first generate region
proposals~\cite{ren2015faster}. They then classify each proposal. This gives
strong accuracy but slower inference. Single-stage detectors remove the
proposal step. They predict boxes directly. YOLO framed detection as one
regression problem~\cite{redmon2016you}. SSD added multi-scale feature
maps~\cite{liu2016ssd}. RetinaNet introduced focal loss for class
imbalance~\cite{lin2017focal}. Later YOLO revisions refined the backbone and
training recipe~\cite{redmon2018yolov3, bochkovskiy2020yolov4}. The
Ultralytics YOLOv8 implementation is our baseline~\cite{jocher2023yolov8}.
EfficientDet made the accuracy--efficiency trade-off
explicit~\cite{tan2020efficientdet}. We adopt that view at the block level for
small-data detection.

\subsection{Attention and Multi-Scale Modules}
Channel attention reweights feature maps by learned importance.
Squeeze-and-Excitation uses a global-pooling squeeze~\cite{hu2018squeeze}. It
then learns per-channel weights with a small network. The parameter overhead
is negligible. CBAM adds spatial attention~\cite{woo2018cbam}. Coordinate
attention encodes position into channel weights~\cite{hou2021coordinate}.
ECA-Net further reduces attention cost with a local one-dimensional
convolution~\cite{wang2020eca}. Multi-scale representation is a complementary
axis. Res2Net builds hierarchical residual connections inside a
block~\cite{gao2021res2net}. Residual learning underpins these deeper
designs~\cite{he2016deep}. Space-to-depth convolution replaces strided
downsampling~\cite{sunkara2022no}. It preserves information for small objects.
Transformer backbones such as ViT~\cite{dosovitskiy2021image} and the
hierarchical Swin Transformer~\cite{liu2021swin} offer another route to global
context. These modules help on large benchmarks. Their value on small data is
less clear.

\subsection{Efficient Architectures and Small Data}
Efficient backbones cut parameters and computation. MobileNetV2 uses
inverted residuals~\cite{sandler2018mobilenetv2}. ShuffleNetV2 uses grouped
convolutions and channel shuffle~\cite{ma2018shufflenet}. GhostNet generates
features cheaply with linear operations~\cite{han2020ghostnet}. All trade some
accuracy for deployability. Model complexity interacts with dataset size.
Normalization and large pretraining reduce
overfitting~\cite{ioffe2015batch, deng2009imagenet}. On small task-specific
data, a larger model is not reliably better. This motivates judging each
change against its cost.

\subsection{Detecting Online Gambling Promotion}
Detecting gambling promotion is a visual content-moderation task. It resembles
logo and brand detection in the wild. The target marks are small and varied.
They appear over dynamic video backgrounds. Public datasets for this task are
scarce. We use two recent community datasets to fill this
gap~\cite{dataset_v21_roboflow, dataset_judi3k_kaggle}. The first isolates
gambling brand logos as objects. The second labels whole reel frames as
gambling or not. Together they cover both fine-grained logo detection and
coarse frame-level screening. This dual view is uncommon in prior logo-detection
work. It lets us test whether a modification helps on the hard case without
hurting the easy one.

\section{Materials and Methods}

\subsection{Datasets}
We use two object-detection datasets, summarized in Table~\ref{tab:datasets}.
Figure~\ref{fig:samples} shows examples from each. The first dataset is
denoted \textbf{Judol Detection v2}. It is sourced from Roboflow
Universe~\cite{dataset_v21_roboflow}. It contains 443 images across five
classes. The classes are online gambling brand logos. They are BK8,
Gate-of-olympus, Princess, Starlight-Princess, and Zeus. Each object is a
compact logo inside a larger scene. This makes Judol Detection v2 a fine-grained detection
task. It is small enough to expose architectural differences. We use it as
the discriminative test bed.

The second dataset is the \textbf{Online Gambling Logo in Instagram Reels
Video} dataset~\cite{dataset_judi3k_kaggle}. We refer to it as \textbf{Instagram
Reels} for brevity. It is collected from Instagram reel videos. It contains 3000
images across two classes. The classes are gambling and nongambling. A key difference matters
here. Each Instagram Reels sample is a full video frame, not a cropped object. The
label describes the whole image. Gambling frames carry promotional overlays
or logos. Nongambling frames show ordinary advertising or lifestyle content.
This is visible in Figure~\ref{fig:samples}(b). The task is closer to
whole-frame classification with localization.

The two datasets play different roles. On Judol Detection v2, models differ by measurable
margins. On Instagram Reels, all evaluated models reach near-ceiling accuracy. We
therefore treat Instagram Reels as a performance-preservation check. It confirms that
no variant regresses on the easier task. Exact licensing follows each source;
Judol Detection v2 is released under CC BY 4.0.

Table~\ref{tab:classdist} lists the Judol Detection v2 class distribution. The five classes
are moderately balanced. Zeus is most common with 229 instances. Starlight-Princess
is least common with 150. BK8 is the second most common class. Yet BK8 is the
hardest class to detect, as shown later. Its difficulty comes from small,
text-like logos, not from scarcity. This distinction guides our error analysis.

\begin{table}[t]
\caption{Characteristics of the two evaluation datasets.}
\label{tab:datasets}
\centering
\footnotesize
\renewcommand{\arraystretch}{1.15}
\begin{tabular}{@{}lccl@{}}
\toprule
\textbf{Dataset} & \textbf{Images} & \textbf{Classes} & \textbf{Evaluation role} \\
\midrule
Judol Detection v2 & 443 & 5 & Architecture discrimination \\
Instagram Reels & 3{,}000 & 2 & Performance preservation \\
\bottomrule
\end{tabular}
\end{table}

\begin{table}[t]
\caption{Class distribution of the Judol Detection v2 dataset (443 images, 932 instances). BK8 is common yet remains the hardest to detect.}
\label{tab:classdist}
\centering
\footnotesize
\renewcommand{\arraystretch}{1.2}
\begin{tabular}{@{}lcc@{}}
\toprule
\textbf{Class} & \textbf{Instances} & \textbf{Share} \\
\midrule
BK8                & 203 & 21.8\% \\
Gate-of-olympus    & 171 & 18.3\% \\
Princess           & 179 & 19.2\% \\
Starlight-Princess & 150 & 16.1\% \\
Zeus               & 229 & 24.6\% \\
\midrule
Total              & 932 & 100\% \\
\bottomrule
\end{tabular}
\end{table}


\begin{figure*}[t]
\centering
\begin{minipage}[b]{0.49\linewidth}
\centering
\includegraphics[width=0.94\linewidth]{figures/samples_v21.png}\\[2pt]
{\footnotesize (a) Judol Detection v2: gambling brand logos in larger scenes.}
\end{minipage}\hfill
\begin{minipage}[b]{0.49\linewidth}
\centering
\includegraphics[width=0.94\linewidth]{figures/samples_judi3k.png}\\[2pt]
{\footnotesize (b) Instagram Reels: full reel frames, gambling (top), nongambling (bottom).}
\end{minipage}
\caption{Dataset examples. Judol Detection v2 targets small logos; Instagram Reels labels whole frames.}
\label{fig:samples}
\end{figure*}

\subsection{Baseline and Variants}
All models share the YOLOv8n backbone, neck, and detection
head~\cite{jocher2023yolov8}. We chose four models to span the complexity axis.
Each represents one clear design strategy. This keeps the comparison
interpretable rather than exhaustive.

\textbf{YOLOv8n (baseline)} is the unmodified nano detector. It fixes the
reference point for accuracy and cost. Every other model is measured against
it. Its 3.16M parameters set the scale for the study.

\textbf{YOLOv8-SE (proposed operating point)} adds squeeze-and-excitation
gates~\cite{hu2018squeeze}. It is a clean ablation of channel attention. Its
backbone and neck are identical to stock YOLOv8. One SE block is inserted on
each detection scale. The blocks sit on the three neck outputs P3, P4, and P5.
Each block acts on its feature map just before the Detect head. No other layer
is changed. Any accuracy difference is therefore attributable to the SE gates
alone.

\textbf{YOLOv8-Res2Net-SPD-SSGA} is the heavy variant. It replaces selected
blocks with Res2Net multi-scale units~\cite{gao2021res2net}. These units widen
the receptive-field range inside a block. It also adds space-to-depth
downsampling~\cite{sunkara2022no}. This operator keeps fine detail that strided
convolution discards. The combination enlarges the model to 4.45M parameters.
It represents the strategy of adding representational capacity.

\textbf{YOLOv8-Shuffle-SE} is the lightweight variant. It is built on
ShuffleNet units with channel shuffle and grouped
convolutions~\cite{ma2018shufflenet}. It keeps a small attention gate for
balance. The design pushes the model below the baseline to 2.70M parameters. It
represents the opposite strategy of trimming capacity for efficiency.

\subsection{Squeeze-and-Excitation Module}
The SE block recalibrates a feature map $U \in \mathbb{R}^{H \times W \times
C}$. It works in two steps. The squeeze step pools spatial information into a
channel descriptor. This gives $z \in \mathbb{R}^{C}$ by global average
pooling:
\begin{equation}
z_c = \frac{1}{H \times W} \sum_{i=1}^{H} \sum_{j=1}^{W} U_c(i,j),
\qquad c = 1,\dots,C .
\end{equation}
The excitation step learns per-channel weights $s \in \mathbb{R}^{C}$. It uses
two fully connected layers with reduction ratio $r$:
\begin{equation}
s = \sigma\!\big(W_2\,\delta(W_1 z)\big),
\end{equation}
where $\delta$ is ReLU and $\sigma$ is the sigmoid. Here $W_1 \in
\mathbb{R}^{\frac{C}{r} \times C}$ and $W_2 \in \mathbb{R}^{C \times
\frac{C}{r}}$. The output rescales each channel of $U$:
\begin{equation}
\tilde{U}_c = s_c \, U_c .
\end{equation}
Only the two small projection matrices are added. The overhead is $2C^2/r$
per block. This is negligible against the backbone. The measured increase is
0.36\% (Table~\ref{tab:relative}). Figure~\ref{fig:se} shows the operation.

The reduction ratio $r$ controls the size of the excitation bottleneck. A
larger $r$ shrinks the two projection matrices further. It lowers the parameter
cost but limits the gate capacity. We keep the standard value from the original
design~\cite{hu2018squeeze}. This choice is deliberate. We aim to test channel
attention as a drop-in module, not to tune it. Placing the gates on the three
neck outputs also matters. Those tensors feed the detection head directly.
Recalibrating them steers the final predictions with minimal disturbance to the
shared backbone.

\subsection{Evaluation Metrics}
We report accuracy with mean average precision. Precision and recall come
from the confusion matrix. Let $\mathrm{TP}$, $\mathrm{FP}$, and
$\mathrm{FN}$ denote true positives, false positives, and false negatives:
\begin{equation}
\mathrm{Precision} = \frac{\mathrm{TP}}{\mathrm{TP}+\mathrm{FP}},
\qquad
\mathrm{Recall} = \frac{\mathrm{TP}}{\mathrm{TP}+\mathrm{FN}} .
\end{equation}
Average precision (AP) is the area under the precision--recall curve. Mean
average precision averages AP over all $K$ classes:
\begin{equation}
\mathrm{AP} = \int_0^1 p(r)\,dr,
\qquad
\mathrm{mAP} = \frac{1}{K}\sum_{k=1}^{K}\mathrm{AP}_k .
\end{equation}
We report mAP50 at IoU 0.5. We also report mAP50-95 averaged over IoU
thresholds 0.5 to 0.95. The normalized confusion matrix entry $C_{ij}$ is the
fraction of true class $i$ predicted as class $j$. Diagonal entries are
correct predictions. Off-diagonal entries are confusions.

\subsection{Fold Construction}
We build five folds with stratified sampling. Cross-validation gives a more
reliable estimate than a single split on small data~\cite{kohavi1995study}. The
strata are the dominant class of each image. Stratification keeps class ratios
stable across folds. It avoids folds that miss a rare class. The split uses a
fixed seed of 0. The same folds
are reused for every model. This removes split luck as a confounder. Each Judol Detection v2
fold holds about 354 training and 89 validation images.

\subsection{Training and Evaluation Protocol}
Each model uses five-fold cross-validation on each dataset. Folds are fixed
across models. Every model sees identical train and validation splits. We
report the mean over the five folds. We also report the fold standard
deviation on Judol Detection v2. This measures run-to-run variability. Model complexity is
parameter count and GFLOPs for one forward pass. It is measured once per
architecture with the Ultralytics profiler.

Training settings are identical across all four models. Each model trains for
100 epochs with early stopping at patience 30. The input size is 640 and the
batch size is 16. The optimizer is SGD with lr0 0.01 and momentum 0.937.
Weight decay is 0.0005 and the seed is fixed at 0. Augmentation uses mosaic,
translation 0.1, scaling 0.5, and horizontal flip 0.5. Mosaic is disabled for
the final ten epochs to stabilize convergence.

We keep the software stack fixed for reproducibility. All runs use a single
NVIDIA RTX A4000 GPU with PyTorch 2.12 and CUDA 13. Deterministic mode is
enabled where supported. Validation uses the best checkpoint of each run.
Complexity is profiled at the same 640 input size. The full workflow is shown
in Figure~\ref{fig:workflow}.

\begin{figure*}[t]
\centering
\includegraphics[width=0.98\linewidth]{fig2.png}
\caption{Evaluation workflow. Both datasets use identical fixed folds and protocol across all four variants.}
\label{fig:workflow}
\end{figure*}

\begin{figure}[t]
\centering
\includegraphics[width=0.60\linewidth]{figures/fig2_se_module.pdf}
\caption{Squeeze-and-excitation channel recalibration.}
\label{fig:se}
\end{figure}

\section{Results and Discussion}

\subsection{Main Comparison}
Table~\ref{tab:main} reports accuracy and complexity for all four models. It
covers both datasets. On the discriminative Judol Detection v2 dataset, YOLOv8-SE leads. It
reaches the best mean mAP50 of 0.8751. It also reaches the best mean mAP50-95
of 0.5896. The baseline follows closely at 0.8736 and 0.5857. The heavier and
lighter variants both fall below the baseline in mAP50-95.

On Instagram Reels, all models cluster tightly. They differ by at most about 0.001
mAP. This confirms that Instagram Reels is saturated. It does not separate the
architectures. We therefore base the architectural discussion on Judol Detection v2. We use
Instagram Reels only to confirm no regression on the easy task.

\begin{table*}[t]
\caption{Accuracy and complexity of the four variants on both datasets. Values are fold means; std is shown for Judol Detection v2. Best per column in bold.}
\label{tab:main}
\centering
\footnotesize
\renewcommand{\arraystretch}{1.2}
\begin{tabular}{@{}lccccrr@{}}
\toprule
& \multicolumn{2}{c}{\textbf{Judol Detection v2}} & \multicolumn{2}{c}{\textbf{Instagram Reels}} & & \\
\cmidrule(lr){2-3}\cmidrule(lr){4-5}
\textbf{Model} & \textbf{mAP50} & \textbf{mAP50-95} & \textbf{mAP50} & \textbf{mAP50-95} & \textbf{Params} & \textbf{GFLOPs} \\
\midrule
YOLOv8n (baseline)          & $0.8736 \pm 0.019$ & $0.5857 \pm 0.025$ & $\mathbf{0.9949}$ & $0.9945$ & $3.157$M & $8.86$ \\
\textbf{YOLOv8-SE (proposed)} & $\mathbf{0.8751} \pm 0.025$ & $\mathbf{0.5896} \pm 0.030$ & $0.9948$ & $\mathbf{0.9948}$ & $3.168$M & $8.87$ \\
YOLOv8-Res2Net-SPD-SSGA     & $0.8716 \pm 0.027$ & $0.5855 \pm 0.028$ & $0.9948$ & $0.9948$ & $4.448$M & $10.98$ \\
YOLOv8-Shuffle-SE           & $0.8675 \pm 0.020$ & $0.5750 \pm 0.027$ & $0.9942$ & $0.9929$ & $\mathbf{2.699}$M & $\mathbf{7.50}$ \\
\bottomrule
\end{tabular}
\end{table*}


The complexity columns add the second half of the story. YOLOv8-SE uses
3.168M parameters and 8.87 GFLOPs. This is nearly identical to the baseline.
The Res2Net variant is far heavier at 4.448M parameters. The Shuffle-SE
variant is the lightest at 2.699M parameters. Accuracy and cost must be read
together. The next subsection makes this explicit.

\subsection{Performance Preservation on Instagram Reels}
The Instagram Reels dataset is a preservation check. All four models score above 0.99
mAP50-95. The task is easy because whole frames carry strong cues. Still, the
ranking is informative. Table~\ref{tab:judi3k} lists the full results. The SE
variant reaches the top mAP50-95 of 0.9948. It also reaches the highest recall
of 0.9954.

\begin{table}[t]
\caption{Results on the Instagram Reels dataset (fold means). All models saturate; SE keeps the best recall and the best mAP50-95. Best per column in bold.}
\label{tab:judi3k}
\centering
\footnotesize
\renewcommand{\arraystretch}{1.3}
\setlength{\tabcolsep}{5pt}
\begin{tabular}{@{}lcccc@{}}
\toprule
\textbf{Model} & \textbf{mAP50} & \textbf{mAP50-95} & \textbf{Prec.} & \textbf{Recall} \\
\midrule
Baseline      & $\mathbf{0.9949}$ & $0.9945$ & $0.9829$ & $0.9927$ \\
SE (ours)     & $0.9948$ & $\mathbf{0.9948}$ & $0.9907$ & $\mathbf{0.9954}$ \\
Res2Net-SPD   & $0.9948$ & $\mathbf{0.9948}$ & $0.9903$ & $0.9940$ \\
Shuffle-SE    & $0.9942$ & $0.9929$ & $\mathbf{0.9920}$ & $0.9932$ \\
\bottomrule
\end{tabular}
\end{table}


The baseline has the lowest precision on Instagram Reels at 0.9829. The SE variant lifts
precision to 0.9907. This means fewer false gambling flags on clean frames.
That property matters for a moderation pipeline. A false flag wastes reviewer
time or blocks benign content.

The lightweight Shuffle-SE has the highest precision but the lowest recall. It
misses slightly more gambling frames. On a screening task, missed positives are
the costly errors. The SE variant offers the better balance here too. It keeps
recall highest without hurting precision much. No variant regresses on this
easy task.

\subsection{Accuracy Versus Complexity}
Table~\ref{tab:relative} expresses each variant as a change from the
baseline. Figure~\ref{fig:tradeoff} plots accuracy against complexity. Three
distinct operating points emerge. Each represents a different design choice.

\begin{table}[t]
\caption{Change of each variant versus the baseline on Judol Detection v2. Positive accuracy and negative complexity are preferable.}
\label{tab:relative}
\centering
\footnotesize
\renewcommand{\arraystretch}{1.3}
\setlength{\tabcolsep}{6pt}
\begin{tabular}{@{}lrrr@{}}
\toprule
\textbf{Model} & \textbf{$\Delta$mAP50-95} & \textbf{$\Delta$Params} & \textbf{$\Delta$GFLOPs} \\
\midrule
SE (ours)     & $+0.0039$ & $+0.36\%$  & $+0.11\%$  \\
Res2Net-SPD   & $-0.0002$ & $+40.90\%$ & $+23.93\%$ \\
Shuffle-SE    & $-0.0107$ & $-14.51\%$ & $-15.35\%$ \\
\bottomrule
\end{tabular}
\end{table}


The SE variant is a near-zero-cost gain. It improves mean mAP50-95 by
$+0.0039$. It adds only 0.36\% parameters and 0.11\% GFLOPs. It is the best
accuracy point at essentially no cost. This is the operating point we
recommend. The Res2Net variant is cost without gain. It adds 40.9\%
parameters and 23.9\% GFLOPs. Yet its mean mAP50-95 is 0.0002 below the
baseline. Added capacity does not help under limited data.

The Shuffle-SE variant is a deliberate trade. It cuts parameters by 14.5\%
and GFLOPs by 15.4\%. It costs 0.0107 mAP50-95. It suits cases where compute
is the binding constraint. In Figure~\ref{fig:tradeoff}, the SE point sits at
the upper left. It combines high accuracy with low computation. The Res2Net
point moves right without moving up. This visual pattern captures the paper's
message.

\begin{figure}[t]
\centering
\includegraphics[width=0.92\linewidth]{figures/fig3_accuracy_complexity.pdf}
\caption{Accuracy--complexity trade-off on Judol Detection v2. Upper-left is preferred.}
\label{fig:tradeoff}
\end{figure}

\subsection{Per-Fold Stability}
Table~\ref{tab:perfold} lists mAP50-95 for every fold. The values reveal wide
spread across folds. The baseline ranges from 0.545 to 0.617. The SE variant
ranges from 0.544 to 0.631. Each model shifts by about 0.07 across folds.

The fold spread is larger than the gap between models. The four models differ
by at most 0.015 in mean mAP50-95. Yet a single model varies by 0.07 across
folds. This tells us the dataset split matters more than the architecture.

We therefore read the model ranking with caution. The per-fold overlap is
substantial. The mean still favors SE, but the margin is small. This is why we
make no claim of statistical significance in this study.

\begin{table}[t]
\caption{Per-fold mAP50-95 on Judol Detection v2 for the four variants. Last column is the mean.}
\label{tab:perfold}
\centering\footnotesize\renewcommand{\arraystretch}{1.3}
\setlength{\tabcolsep}{5pt}
\begin{tabular}{@{}lcccccc@{}}
\toprule
\textbf{Model} & \textbf{F1} & \textbf{F2} & \textbf{F3} & \textbf{F4} & \textbf{F5} & \textbf{Mean} \\
\midrule
Baseline & 0.6052 & 0.5827 & 0.5787 & 0.5449 & 0.6169 & 0.5857 \\
SE (ours) & 0.6016 & 0.5719 & 0.6004 & 0.5438 & 0.6305 & 0.5896 \\
Res2Net-SPD & 0.6260 & 0.5714 & 0.5813 & 0.5441 & 0.6046 & 0.5855 \\
Shuffle-SE & 0.6101 & 0.5721 & 0.5848 & 0.5270 & 0.5809 & 0.5750 \\
\bottomrule
\end{tabular}
\end{table}


\subsection{Precision and Recall}
Table~\ref{tab:pr} reports precision and recall on Judol Detection v2. The baseline has the
highest precision at 0.942. The SE variant trades a little precision for a
touch more recall. Its recall is 0.818 against the baseline 0.816.

The lightweight Shuffle-SE has the lowest precision at 0.894. It emits more
false positives than the others. This is consistent with its lower mAP. The
Res2Net variant reaches the highest recall at 0.824.

These differences are small and mostly within fold noise. The main message is
balance. The SE variant keeps precision and recall close to the baseline. It
does not sacrifice one for the other. This stable behavior suits a moderation
setting.

\begin{table}[t]
\caption{Mean precision and recall on Judol Detection v2 over five folds. Best per column in bold.}
\label{tab:pr}
\centering\footnotesize\renewcommand{\arraystretch}{1.3}
\begin{tabular}{@{}lcc@{}}
\toprule
\textbf{Model} & \textbf{Precision} & \textbf{Recall} \\
\midrule
Baseline & \textbf{0.9424}$\pm$0.017 & 0.8157$\pm$0.030 \\
SE (ours) & 0.9318$\pm$0.015 & 0.8178$\pm$0.040 \\
Res2Net-SPD & 0.9335$\pm$0.030 & \textbf{0.8235}$\pm$0.030 \\
Shuffle-SE & 0.8945$\pm$0.033 & 0.8186$\pm$0.028 \\
\bottomrule
\end{tabular}
\end{table}


\subsection{Training Convergence}
Figure~\ref{fig:conv} plots the fold-mean mAP50-95 against the epoch. All four
curves rise sharply in the first twenty epochs. They then flatten and plateau
after about sixty epochs. The four curves stay close throughout training.

The SE variant reaches 0.5 mAP50-95 a few epochs earlier than the baseline. It
crosses the threshold near epoch 37 against epoch 42. This suggests that
channel recalibration slightly speeds early learning. The effect is modest but
consistent.

The lightweight Shuffle-SE curve sits below the others late in training. It
converges to a lower plateau. This matches its lower final accuracy in
Table~\ref{tab:main}. The convergence view confirms the accuracy ranking is
stable, not an artifact of a single checkpoint.

\begin{figure}[t]
\centering
\includegraphics[width=0.86\linewidth]{figures/convergence.pdf}
\caption{Fold-mean mAP50-95 versus epoch on Judol Detection v2 for the four variants.}
\label{fig:conv}
\end{figure}

\subsection{Confusion Matrix Analysis}
Figure~\ref{fig:cm} shows the normalized confusion matrix. It is computed for
YOLOv8-SE on a Judol Detection v2 validation fold. The diagonal is generally strong. Four
of five classes exceed 0.85 recall. Starlight-Princess is the easiest class
at 0.96. Zeus and Gate-of-olympus also detect reliably.

The matrix reveals where errors concentrate. BK8 is the hardest class. Only
0.52 of BK8 instances are detected correctly. About 0.48 are confused with
background. This means many BK8 logos are missed. The background column is the
main error source across classes. This is typical for small logos in cluttered
scenes.

These errors are informative for future work. BK8 logos vary widely in size
and placement. They often appear as small overlays. A detector may treat them
as background. Class-specific augmentation could help here. The matrix also
shows little confusion between logo classes. The classes are visually
distinct. The main challenge is detection, not classification.

The background confusion has a clear moderation cost. A missed BK8 logo is a
gambling promo that slips through. Raising recall on small marks is therefore a
priority. Higher-resolution inference could help recover these logos. Tiling
the input is another option for small-object recall. We leave both directions
to future work.

\begin{figure}[t]
\centering
\includegraphics[width=0.90\linewidth]{figures/confusion_matrix.png}
\caption{Normalized confusion matrix for YOLOv8-SE on a Judol Detection v2 fold.}
\label{fig:cm}
\end{figure}

\subsection{Qualitative Detection Results}
Figure~\ref{fig:det} shows detections from YOLOv8-SE. The images are Judol Detection v2
validation samples. The model localizes gambling logos accurately. Confidence
scores are generally high. Gate-of-olympus is detected at 0.94. Two
Starlight-Princess logos are detected at 0.95.

The examples cover varied backgrounds. Some logos sit on bright game art.
Others sit on sports or advertising scenes. The model handles both cases. It
places tight boxes around each logo. This matches the strong diagonal in the
confusion matrix.

One BK8 example is detected at only 0.69. This lower confidence agrees with
the confusion matrix. BK8 remains the least reliable class. The qualitative
and quantitative views are consistent. Together they build a complete picture
of model behavior.

\begin{figure}[t]
\centering
\includegraphics[width=0.92\linewidth]{figures/detections.png}
\caption{YOLOv8-SE detections on Judol Detection v2 validation images.}
\label{fig:det}
\end{figure}

\subsection{Grad-CAM Interpretability}
We use Grad-CAM to inspect where the model looks. The heatmaps are computed
for YOLOv8-SE on Judol Detection v2 images. We aggregate gradients over the three
SE-recalibrated detection scales P3, P4, and P5. This covers logos of
different sizes. Warm regions drive the strongest class response.
Figure~\ref{fig:gradcam} shows the result.

The heatmaps concentrate on the gambling brand identity. The BK8 wordmark is
clearly highlighted in two examples. The Gates-of-Olympus and Starlight
Princess artwork also draw strong attention. The model focuses on the logo,
not the surrounding scene. This behavior is what a reliable detector should
show.

The maps also explain the SE contribution qualitatively. SE gates reweight
channels before detection. This sharpens the response on discriminative
regions. The focused heatmaps are consistent with the accuracy results. They
show that attention lands on the right visual cues. This supports the use of
channel recalibration for this task.

\begin{figure}[t]
\centering
\includegraphics[width=0.92\linewidth]{figures/gradcam.png}
\caption{Grad-CAM for YOLOv8-SE on Judol Detection v2, aggregated over the P3/P4/P5 scales.}
\label{fig:gradcam}
\end{figure}

\subsection{Instagram Reels: Detection and Interpretability}
We repeat the qualitative analysis on the Instagram Reels dataset. This dataset
labels whole frames, not cropped logos. Figure~\ref{fig:det_judi} shows
YOLOv8-SE predictions on gambling reels. The model flags the frames with high
confidence. It fires on frames that carry promotional overlays and betting
brand marks. This confirms the strong numbers in Table~\ref{tab:judi3k}.

Figure~\ref{fig:cm_judi} shows the confusion matrix for the two classes. Both
classes sit near the diagonal. Gambling reaches 1.00 and nongambling reaches
0.99. The classes almost never swap. A small share of background regions are
predicted as one class or the other. This is expected when a detector is
applied to whole-frame labels. It does not affect the frame-level decision.

Figure~\ref{fig:gc_judi} shows Grad-CAM on the same reels. The pattern differs
from Judol Detection v2. On logos, the heat lands on the brand mark. On reels,
the heat spreads over the central subject and foreground. The model reads the
whole frame, not a single logo. This matches the dataset design. It also
explains why the task saturates: many holistic cues signal a gambling reel.

\begin{figure}[t]
\centering
\includegraphics[width=0.99\linewidth]{figures/detections_judi3k.png}
\caption{YOLOv8-SE detections on gambling Instagram reels (whole-frame labels).}
\label{fig:det_judi}
\end{figure}

\begin{figure}[t]
\centering
\includegraphics[width=0.72\linewidth]{figures/confusion_matrix_judi3k.png}
\caption{Confusion matrix for YOLOv8-SE on the Instagram Reels dataset.}
\label{fig:cm_judi}
\end{figure}

\begin{figure}[t]
\centering
\includegraphics[width=0.99\linewidth]{figures/gradcam_judi3k.png}
\caption{Grad-CAM on Instagram reels. Attention spreads over the whole frame, not one logo.}
\label{fig:gc_judi}
\end{figure}

\subsection{Per-Class Detection Difficulty}
Table~\ref{tab:perclass} and Figure~\ref{fig:perclass} break accuracy down by
class. Four of the five classes are easy. Gate-of-olympus, Princess,
Starlight-Princess, and Zeus all exceed 0.92 mAP50. These logos have large,
distinctive artwork.

BK8 is the hardest class by a wide margin. Its mAP50 falls between 0.56 and
0.61 for every model. This gap dominates the overall differences between
models. The BK8 wordmark is small and text-like. It blends into busy sports
and advertising backgrounds.

No single model wins on BK8 across the board. The baseline and Shuffle-SE
score highest on BK8 in this fold. The SE variant leads on the mean but not on
BK8. This shows the overall ranking is driven by the easy classes. Improving
BK8 is the clearest path to a better detector.

\begin{table*}[t]
\caption{Per-class mAP50 on Judol Detection v2 (fold 0). BK8 is the hardest class for every model.}
\label{tab:perclass}
\centering\footnotesize\renewcommand{\arraystretch}{1.2}
\begin{tabular}{@{}lccccc@{}}
\toprule
\textbf{Model} & \textbf{BK8} & \textbf{Gate} & \textbf{Princess} & \textbf{Starlight} & \textbf{Zeus} \\
\midrule
YOLOv8n (baseline) & 0.612 & 0.941 & 0.969 & 0.995 & 0.934 \\
YOLOv8-SE (proposed) & 0.564 & 0.939 & 0.942 & 0.986 & 0.949 \\
YOLOv8-Res2Net-SPD-SSGA & 0.589 & 0.967 & 0.949 & 0.992 & 0.945 \\
YOLOv8-Shuffle-SE & 0.604 & 0.931 & 0.933 & 0.992 & 0.928 \\
\bottomrule
\end{tabular}
\end{table*}


\begin{figure*}[t]
\centering
\includegraphics[width=0.82\linewidth]{figures/perclass.pdf}
\caption{Per-class mAP50 on Judol Detection v2. BK8 is far harder than the other four classes.}
\label{fig:perclass}
\end{figure*}

\subsection{Broader Architectural Landscape}
We also place SE within a wider pool of nano variants. Table~\ref{tab:broader}
and Figure~\ref{fig:broader} compare SE against these alternatives. The pool
covers many attention, multi-scale, and efficient designs. All use the same
protocol and folds.

SE sits in the top tier of the explored models. Against every alternative in
Table~\ref{tab:broader}, none does better. Several are clearly heavier yet
score lower. They add 20 to 40 percent more parameters for no gain. SE reaches
its accuracy at almost the baseline cost.

This restates the core efficiency argument. SE is the smallest change to the
standard YOLOv8n. It is a drop-in module with well-understood behavior. It
matches or beats far more elaborate redesigns at a fraction of their cost. This
makes it a strong default when engineering effort is limited.

\begin{table}[t]
\caption{Proposed SE model ($\star$) against other nano-scale variants on Judol Detection v2, with cost. None of these simpler or heavier alternatives beats it.}
\label{tab:broader}
\centering\footnotesize\renewcommand{\arraystretch}{1.2}
\begin{tabular}{@{}lccc@{}}
\toprule
\textbf{Model} & \textbf{mAP50-95} & \textbf{Params (M)} & \textbf{GFLOPs} \\
\midrule
yolov8-se-baseline $\star$ & 0.5896 & 3.17 & 8.9 \\
yolov8-mbconv-coordatt-simam & 0.5887 & 3.10 & 8.1 \\
yolov8-swin-ema & 0.5872 & 3.01 & 20.8 \\
yolo12-spd-shuffle & 0.5865 & 4.62 & -- \\
yolov8-mbconv-eca & 0.5858 & 3.09 & 8.1 \\
yolov8n-base (base) & 0.5857 & 3.16 & 8.9 \\
yolov8-res2net-spd-ssga & 0.5855 & 4.45 & 11.0 \\
yolov8-ghostv2-spd-gam & 0.5841 & 5.66 & 11.8 \\
\bottomrule
\end{tabular}
\end{table}


\begin{figure}[t]
\centering
\includegraphics[width=0.86\linewidth]{figures/broader.pdf}
\caption{Accuracy versus parameters for top nano variants. SE is near-top at low cost.}
\label{fig:broader}
\end{figure}

\subsection{Discussion}
The pattern across the three modifications is consistent. In a small-data
regime, added capacity is a liability. The two capacity-changing variants
both underperform the baseline. Enlargement and reduction both hurt mean
accuracy. Only the SE variant improves it. SE changes almost no capacity. It
recalibrates existing features instead. This aligns with the low-overhead
design of channel attention~\cite{hu2018squeeze}.

We stress the modest size of the SE gain. The improvement on Judol Detection v2 is small in
absolute terms. Its value lies in cost, not magnitude. It achieves the best
accuracy at the lowest incremental cost. This is exactly the
accuracy--efficiency argument of the paper. The message is a design
principle, not a leaderboard claim.

The result has practical weight for moderation deployment. A gambling detector
may run on millions of frames per day. Every extra GFLOP raises the compute
bill at that scale. The SE variant adds only 0.11\% GFLOPs over the baseline.
It therefore improves accuracy essentially for free. A heavier redesign would
raise cost with no reliable accuracy return here. Efficiency-first tuning is
the rational choice for this setting.

The broader benchmark supports the same view from a wider angle. Across many
nano variants, most complex redesigns cluster near the baseline. Only a handful
edge ahead, and never by a large margin. This ceiling likely reflects the
dataset, not the architectures. When data is the bottleneck, architecture
tuning has diminishing returns. Investment in more and better labeled data may
help more than a heavier model.

\subsection{Threats to Validity}
We group the limitations by the kind of validity they affect.

\textit{Statistical conclusion validity.} The discriminative dataset holds
only 443 images. The comparison uses five folds. The fold spread exceeds the
gap between models (Table~\ref{tab:perfold}). We therefore report the fold
standard deviation and make no significance claim. The ranking is indicative,
not proven.

\textit{Internal validity.} All models share one training protocol, one seed,
and identical folds. This isolates the architecture as the only variable. A
single seed still leaves run-to-run noise unmeasured. Multiple seeds would
tighten the estimates. The per-class table uses one fold and should be read as
illustrative.

\textit{External validity.} The two datasets come from specific sources. One
is a logo set; the other is Instagram reel frames. Both center on Indonesian
online gambling promotion. Results may not transfer to other platforms or
regions. The saturation on Instagram Reels also limits what it can tell us.

\textit{Construct validity.} We measure accuracy with mAP and cost with
parameters and GFLOPs. GFLOPs is a proxy for latency, not a direct measure.
Real deployment speed depends on hardware and the runtime. We report GFLOPs as
a comparable proxy within this study only.

\section{Conclusion}
We compared three architectural modifications of YOLOv8n. They are channel
attention, multi-scale enlargement, and lightweight reduction. We judged each
by accuracy and by the complexity it costs. We used two datasets of
contrasting size. On the small discriminative dataset, SE gave the best mean
accuracy. It added only 0.36\% parameters. Structural enlargement added 40.9\%
parameters without any gain. Lightweight reduction traded accuracy for a
14.5\% parameter saving. A broader benchmark of nano variants placed SE in the
top tier at near-baseline cost. Per-fold, per-class, and convergence analyses
all agreed with this ranking.

The practical takeaway is simple. Prefer cheap feature recalibration over
structural growth on small data. Report accuracy jointly with its cost. For
gambling-content moderation, this keeps compute low at scale while holding
accuracy. The hardest remaining case is the small, text-like BK8 logo, which
limits every model. Future work should target that case directly. It should
extend the fold count and dataset diversity to enable formal significance
testing. It should study SE insertion placement and reduction ratio
systematically. It should also measure real inference latency on target
hardware, not GFLOPs alone.

\section*{Acknowledgment}
Acknowledgments, funding details, and the names of contributing participants
are withheld to preserve anonymity for double-blind review. They will be
restored in the camera-ready version.

\bibliographystyle{IEEEtran}
\bibliography{references}

\end{document}
